\documentclass{alex_summary}

\begin{document}
\section*{Reading summary}

As the title suggests, \blockcquote{devoret_superconducting_2013}{Superconducting Circuits for Quantum Information: An Outlook} gives a superficial perspective on quantum computation using the superconducting circuits platform. They authors explain past efforts leading up to today and then go on to describe how they envision the how future research in this area will carry out. \blockcquote{gao_coherent_2015}{Coherent manipulation, measurement and entanglement of individual solid-state spins using optical fields}

Superconducting circuits as possible platform for quantum computation have emerged in the late 1990s and have since been at the forefront of modern research. They leverage \textbf{superconductivity} and the \textbf{Josephson effect} to create \textbf{anharmonic LC oscillators} with quality factors in the range of \( \mathbf{10^6} \) resulting in lifetimes of \( \sim \mathbf{100 \unit{\micro\s}} \). State of the art lithography pioneered by the semiconductor industry is used to pattern the electrical circuit onto a wafer to create capacitive and inductive elements (forming the qubit), that are connected via transmission lines. \textbf{Microwave signals} can then be used for qubit manipulation and readout. 

The outstanding feature of superconducting circuits (for me) lies in their flexibility. Whilst most other platforms rely on \blockquote{naturally occurring} two-level systems, the two-level structure of superconducting qubits is an emergent property of the electrical circuit design (which is why they are often referred to as artificial atoms). This has both advantages and drawbacks. On the one hand, we can tune parameters such as qubit frequency or coupling strength to suit ones need. On the other hand, being to manipulate such fundamental properties implies that . For this reason superconducting qubits need to be characterized frequently to account for drift

\begin{question1}

\end{question1}


\begin{question3}

\end{question3}

\printbibliography
\end{document}
